%% kapa: exam=2 round=27 year=2016 period=1 subject=appraisal_practice subject_ko=감정평가실무 seq=2 number=2 points=30
%% source: Q-Net 감정평가사 기출문제 (cst00309) artlSeq=5185401 file=extracted/1008/5185401/제27회 감정평가사 2차시험 기출문제/제27회 감정평가사 2차 시험문제지/(1교시 시험문제지)감정평가실무.pdf
%% conversion: kordoc 4.24.51 -> tools/gen2.py (auto); review=proofread-gpt-sol; flags=none
\begin{longquestion}{【문제 2】}
감정평가사 甲은 임대에 제공되고 있는 상업용 부동산(집합건물)을 시장가치로 매수할 것을 제안받은 잠재적 매수인 乙로부터 상담을 의뢰받았는 바, 제시된 자료를 참조하여 다음 물음에 답하시오. (30점) \par
\medskip\noindent 물음 1) 의뢰인 乙이 제안을 받아들일 경우, 乙의 투자수익률(거래비용은 고려하지 않음)과 乙의 요구수익률을 충족시키는 매매가격을 구하시오. (10점) \par
\medskip\noindent 물음 2) 본건 부동산의 완전소유권에 기초한 수익률을 산출하여 직접환원법에 의한 수익가치를 구하시오. (10점) \par
\medskip\noindent 물음 3) 「완전소유권의 가치 = 임대권의 가치 + 임차권의 가치」라는 등식이 성립하고, 의뢰인 乙의 요구수익률을 충족시키는 매매가격이 적정한 임대권의 가치라 가정한다. 이 경우, 내재된 임차권 수익률을 구하고, 임차권 수익률이 임대권 수익률보다 큰 이유를 설명하시오. (10점) \par
<자료 1> 기본적 사항 \par
1. 의뢰일(기준시점) : 2016년 7월 1일 \par
2. 대상부동산 내역 \par
{\small\centering
\resizebox{\ifdim\width>\linewidth\linewidth\else\width\fi}{!}{%
\begin{tabular}{|c|c|c|c|c|c|}
\hline
\multirow{2}{*}{\makecell{소재지 지번}} & \multirow{2}{*}{\makecell{층ㆍ호수}} & \multirow{2}{*}{\makecell{면적\\(m\textsuperscript{2})}} & \multicolumn{2}{|c|}{\makecell{임대료(원)}} & \multirow{2}{*}{\makecell{임대차\\계약기간}} \\
\cline{4-5}
 &  &  & \makecell{보증금} & \makecell{지불임료/월} &  \\
\hline
\makecell{S시 S구 S동\\1000번지} & \makecell{지상2층\\201호} & \makecell{100.0} & \makecell{100,000,000} & \makecell{1,500,000} & \makecell{2015.07.01\\～ 2020.06.30} \\
\hline
\end{tabular}}\par}

※ 계약내용이나 조건은 전형적인 시장 관행에 따른 것임 \par
<자료 2> 의뢰인 면담 내용 요약 \par
1. 무위험률에 1.20 \% 가산한 수익률을 요구하며, 이러한 조건은 과거 상업용 부동산(집합건물) 투자수익률과 무위험률의 격차(spread) 추세 및 인근지역의 전형적인 임대용 부동산 투자자의 요구 조건과도 부합함 \par
2. 본건에 투자할 경우, 임대차 계약이 만료되는 시점에 재매도 예정임 \par
<자료 3> 시장조사 자료 \par
1. 시장임대료(의뢰일 현재) : 보증금 @1,100,000원/m\textsuperscript{2} 및 지불임료 @16,500원/m\textsuperscript{2} 수준이고, 관리비는 별도로 임차인이 지불함 \par
2. 매매사례 \par
{\small\centering
\resizebox{\ifdim\width>\linewidth\linewidth\else\width\fi}{!}{%
\begin{tabular}{|c|c|c|c|c|c|c|}
\hline
\makecell{기호} & \makecell{소재지\\지번} & \makecell{층ㆍ호수} & \makecell{면적\\(m\textsuperscript{2})} & \makecell{계약일} & \makecell{매매금액(원)} & \makecell{비고} \\
\hline
\makecell{1} & \makecell{S시 S구\\S동\\986번지} & \makecell{지상2층\\201호} & \makecell{30.0} & \makecell{2016.06.24} & \makecell{363,000,000} & \makecell{경사지에\\위치한바, 주도로기준 1층} \\
\hline
\makecell{2} & \makecell{S시 S구\\S동\\1021번지} & \makecell{지상2층\\205호} & \makecell{210.0} & \makecell{2016.06.19} & \makecell{1,000,000,000} & \makecell{-} \\
\hline
\makecell{3} & \makecell{S시 S구\\S동\\1004번지} & \makecell{지상2층\\207호} & \makecell{92.0} & \makecell{2016.06.20} & \makecell{540,000,000} & \makecell{-} \\
\hline
\end{tabular}}\par}

※ 상기 매매사례 모두 최근의 것으로 별도의 시점수정은 불필요함 \par
3. 개별요인 \par
{\small\centering
\resizebox{\ifdim\width>\linewidth\linewidth\else\width\fi}{!}{%
\begin{tabular}{|c|c|c|c|}
\hline
\makecell{본건} & \makecell{매매사례 \#1} & \makecell{매매사례 \#2} & \makecell{매매사례 \#3} \\
\hline
\makecell{1.00} & \makecell{0.50} & \makecell{0.85} & \makecell{1.03} \\
\hline
\end{tabular}}\par}

4. 본건 부동산의 인근지역 내 상업용 부동산(집합건물)의 공급 증가로 본건 임대차계약 만료 시에 5.00 \%의 시장가치 하락이 예상됨 \par
<자료 4> 시장금리 등 자료 \par
1. 국고채(3년) : 1.390 \%/년 (최근 3개월 평균) \par
2. 회사채(장외3년, AA- 등급) : 1.943 \%/년 (최근 3개월 평균) \par
3. 보증금 운용이율 : 2.00 \%/년 \par
\end{longquestion}
